\documentclass[twoside]{report}
\usepackage[utf8]{inputenc}
\usepackage[greek,english]{babel}
\usepackage{graphicx}
\usepackage{float}
\usepackage{geometry}
\usepackage{lipsum}
\usepackage{minted}
\usepackage{xcolor} %to access the named color LightGray in minted environment
\usepackage{fancyhdr}
\usepackage{setspace}
\usepackage{amsmath}
\usepackage[T1]{fontenc}
\usepackage{mathpazo}
\usepackage{afterpage}
\usepackage[intoc]{nomencl}
\usepackage{etoolbox}

\makenomenclature

\newcommand\blankpage{%
    \null
    \thispagestyle{empty}%
    \addtocounter{page}{}%
    \newpage}

\usepackage{tikz}
\usetikzlibrary{calc}
\newcommand\HRule{\rule{\textwidth}{1pt}}

\pagestyle{fancy}
\fancyhf{}
\rhead{Alexandros Dimas}
\lhead{Motion Planning for Collision-tolerant Aerial Robots}
\lfoot{\small Department of Mechanical Engineering & Aeronautics - Division of Design and Manufacturing}
\renewcommand{\footrulewidth}{1pt}
\rfoot{\thepage}
\cfoot{}
\makeatletter
\renewcommand\chapter{\if@openright\cleardoublepage\else\clearpage\fi
                    \thispagestyle{fancy}%
                    \global\@topnum\z@
                    \@afterindentfalse
                    \secdef\@chapter\@schapter}
\makeatother


% \newgeometry{
%     top=1.25in,
%     bottom=1.25in,
%     outer=0.75in,
%     inner=1.5in,
% }

\title{Motion Planning for Collision-tolerant Aerial Robots}
\author{Alexandros Dimas}
\date{\today}
\renewcommand{\baselinestretch}{1.5} 

\begin{document}

\begin{titlepage}

\begin{tikzpicture}[remember picture, overlay]
  \fill [lightgray] ($(current page.north west) + (1in,-1in)$) rectangle ($(current page.south east) + (-1in,1in)$);
\end{tikzpicture}

\begin{minipage}{0.25\textwidth}
\begin{figure}[H]
\includegraphics[width =\textwidth]{images/logo-up-4color-stamp.png}
\end{figure}
\end{minipage} \hfill
\begin{minipage}{0.85\textwidth}
{\setstretch{2}
\large \textbf{University of Patras}\\
\large \textbf{Department of Mechanical Engineering and Aeronautics}\\
\large \textbf{Division of Design and Manufacturing}\\
\large \textbf{Robotics Group}
}
\end{minipage}

\vspace{1cm}

\centering
\large{\textbf{Diploma Thesis}}

\vspace{1cm}

\Large{\textbf{Motion Planning for Collision-Tolerant Aerial Robots}}

\vspace{1.2cm}

\Large{\textbf{Alexandros Dimas}}\\
\large 1054531

\vspace{1.2cm}
\textbf{Evagellos Dermatas, Associate professor, University of Patras}\\
\textbf{Panagiotis Koustoumpardis,\ Research \& Teaching Associate , \\University of Patras}\\
\textbf{Kostas Alexis,\ Full Professor, Norwegian University of Science and Technology}

\vspace{1.2cm}
dissertation submitted to the department of Mechanical Engineering and Aeronautics of University of Patras

\vspace{0.2cm}
\Large Patras, Month, 2023
\end{titlepage}
\newpage

\pagenumbering{roman}

\mbox{}
\vfill

\noindent
{\tt
University of Patras, Department of Mechanical Engineering and Aeronautics\\
Alexandros Dimas\\ \copyright\ 2023 -- All rights reserved \newline
}
\newpage
\begin{minipage}{0.25\textwidth}
\begin{figure}[H]
\includegraphics[width =\textwidth]{images/logo-up-4color-stamp.png}
\end{figure}
\end{minipage} \hfill
\begin{minipage}{0.85\textwidth}
{\setstretch{2}
\large \textbf{University of Patras}\\
\large \textbf{Department of Mechanical Engineering and Aeronautics}\\
\large \textbf{Division of Design and Manufacturing}\\
\large \textbf{Robotics Group}
}
\end{minipage}

\vspace{4cm}
\begin{center}
    This dissertation was presented \\ by \textbf{Alexandros Dimas} \\ 1054531 \\ 
    \vspace{2cm}
    on \today
\end{center}

\newpage
\mbox{}
\vfill
\noindent
\textgreek{
Η έγκριση της διπλωματικής εργασίας δεν υποδηλοί την αποδοχή των γνωμών του συγγραφέα.\\ Κατά τη συγγραφή τηρήθηκαν οι αρχές της ακαδημαϊκής δεοντολογίας.} 
{\tt
}
\newpage

\clearpage
\addcontentsline{toc}{chapter}{\textgreek{Περίληψη}}
\chapter*{\textgreek{Περίληψη}}
\textgreek{
\begin{center}
    \textbf{
    \large Σχεδιασμός Κίνησης για Εναέρια Ρομποτικά Οχήματα με Ανοχή σε Συγκρούσεις
    \\
    \large Αλέξανδρος Δήμας
    }
\end{center}
[Εδώ μπαίνει η περίληψη της εργασίας που δεν θα πρέπει να υπερβαίνει τις 350 λέξεις]
}
\mbox{}
\vfill
\noindent
\textgreek{
Λέξεις κλειδιά \\ 
Εδώ μπαίνουν μέχρι 5 λέξεις κλειδιά που αντιστοιχούν στα πλέον αντιπροσωπευτικά αντικείμενα-θέματα της εργασίας  
}
\afterpage{\null\newpage}
\newpage

\clearpage
\addcontentsline{toc}{chapter}{Abstract}
\chapter*{Abstract}
\begin{center}
    \textbf{
    \large Motion Planning for Collision-Tolerant Aerial Robots
    \\
    \large Alexandros Dimas
    }
\end{center}
\lipsum[1-2]
\mbox{}
\vfill
\noindent
Keywords \\ Up to five relevant keywords
\afterpage{\null\newpage}

\clearpage
\addcontentsline{toc}{chapter}{List of Tables}
\listoftables
\afterpage{\null\newpage}
\clearpage
\addcontentsline{toc}{chapter}{List of Figures}
\listoffigures
\afterpage{\null\newpage}

% https://www.overleaf.com/learn/latex/Nomenclatures
\newpage
\chapter*{Nomenclature}
\addcontentsline{toc}{chapter}{Nomenclature}

\section*{Abbreviations}

\begin{table}[h]

\label{NOMENCLATURE}

\begin{tabular}{ll}

	\hline

	\textbf{Abbreviation} & \textbf{Definition}	\\ \hline

	aT          	& a Thing  			\\

	aST		& a Second Thing 		\\ \hline

\end{tabular}

\end{table}
\section*{Symbols}


\begin{tabular}{ll}

	\hline

	\textbf{Symbol} & \textbf{Unit}	\\ \hline

	aT          	& a Thing  			\\

	aST		& a Second Thing 		\\ \hline

\end{tabular}
\afterpage{\null\newpage}


\tableofcontents
\afterpage{\null\newpage}

\clearpage
\pagenumbering{arabic}
\addcontentsline{toc}{chapter}{Acknowledgements}
\chapter*{Acknowledgements}
I want to thank ...
\lipsum[1-2]
\afterpage{\null\newpage}


\chapter{Introduction}
\lipsum[1-2]
\section{Related Work}
\lipsum[1-10]
\section{Motivation and Contributions}
\lipsum[1-10]
\chapter{Modelling and Control}

\section{The mechanics of a quadrotor}
This section will explain how the quadrotor moves in each direction as well as how it is able to rotate around each of the body-fixed axes.
\lipsum[1-10]
\section{Deriving the equations of motion}
The Newton-Euler derivation of the equation of motion.

\begin{equation}
    m \ddot{\textbf{r}} = \begin{bmatrix}
        0 \\ 0 \\ -mg
    \end{bmatrix}
    + R \begin{bmatrix}
    0 \\ 0 \\ F_1+F_2+F_3+F_4
    \end{bmatrix}
\end{equation}

\begin{equation}
    I \begin{bmatrix} \dot{p} \\ \dot{q} \\ \dot{r} \end{bmatrix} = \begin{bmatrix} L(F_2-F_4) \\ L(F_3-F_1) \\ M_1-M_2+M_3-M_4 \end{bmatrix} - \begin{bmatrix} p \\ q \\ r \end{bmatrix} \times I \begin{bmatrix} p \\ q \\ r \end{bmatrix}
\end{equation}

\lipsum[1-10]
\begin{figure}
    \centering
    \includegraphics[width=\textwidth]{images/4464460_orig.png}
    \caption{The world and body-fixed frames}
    \label{fig:frames}
\end{figure}
\chapter{Motion Planning}

I can reference LaValle \cite{LaValle2006}
\lipsum[1-10]

\clearpage
\addcontentsline{toc}{chapter}{Bibliography}
\bibliographystyle{plain}
\bibliography{references}

\appendix
\chapter{C++ Code}
\begin{minted}
[
frame=single,
framesep=2mm,
baselinestretch=1.2,
bgcolor=lightgray,
fontsize=\footnotesize,
linenos
]
{cpp}
#include <iostream>

int main()
{
    std::cout << "Hello World" << '\n';
    return 0;
}
\end{minted}
\include{matlab-code}
\end{document}